% ================================================================
% ==== FILE: content/k_levels/k12.tex
% ================================================================

\paragraph{Canonical tuple projection note.}
Local K-level displays in this module are projections of the canonical public
tuple \(K=(\Omega,\partial\Omega,A,\Theta,P,J,C,k,M)\). When a display omits
\(M\), reorders components, or uses level-indexed shorthand, the omitted
embedding context is inherited from the surrounding level discussion rather
than introduced as a competing definition.

\subsubsection{\texorpdfstring{$K_{12}$}{K12} as a Provisional Upper-Coherence Candidate}
\label{sec:k12-overview}

\(K_{12}\) is retained in Core~1.3.3 as an upper-taxonomy research candidate.
It is not promoted as a completed limit theorem, a final level of reality, or a
proof that all branches, flows, operators, or time directions have globally
closed. The public role of the module is to state what an upper-coherence claim
would have to prove before stronger wording could be licensed.

A cautious candidate state description is
\[
\Omega(K_{12}) =
\{\Omega^*,A^*,P^*,\Theta^*,J^*,C^*,\mu^{(12)}\},
\]
where the starred objects denote proposed upper-level projections of the lower
K-level vocabulary. In this release they are research objects. They become
support-carrying only when a later proof or evidence package names the
operator assumptions, the convergence criterion, the boundary condition, the
negative case, and the demotion rule.

\subsubsection{Candidate Components and Required Proof Obligations}

\begin{itemize}
\item \(\Omega^*\) would have to identify admissible upper-coherence states
      without absorbing incompatible \(K_0\)--\(K_{11}\) witnesses by rhetoric.
\item \(A^*\) would have to show which upper-level axes are retained and which
      apparent axes are demoted as inert.
\item \(\Theta^*\) would have to state threshold conditions under which an
      upper-coherence reading is allowed or rejected.
\item \(P^*\) and \(J^*\) would have to specify potentials and flows under
      explicit operator assumptions rather than by global-invariance language.
\item \(C^*\) would have to identify cycles that support the candidate reading
      and cases where such cycles fail.
\item \(M_{12}\) would have to name the embedding context that makes the
      candidate admissible.
\end{itemize}

These requirements are intentionally strict. They prevent the upper taxonomy
from turning into a totalizing closure claim.

\subsubsection{Promotion Checklist}

An upper-coherence statement may be promoted only after it passes a checklist
that is stronger than ordinary exposition. The statement must identify the
candidate object, the lower-level witnesses retained by projection, the
operator assumptions, the convergence or non-convergence criterion, the
boundary condition, the finite or formal witness if one exists, the comparator
that could explain the same structure with less burden, and the condition under
which the claim is demoted.

\begin{center}
\begin{tabular}{p{0.25\linewidth}p{0.62\linewidth}}
\toprule
Review item & Required public answer \\
\midrule
Object & What is the specific \(K_{12}\) candidate, and which lower-level
witnesses does it retain? \\
Assumptions & Which operator, embedding, and projection assumptions are being
used? \\
Support & Is the statement definitional, proof-bound, finite-witness-bound, or
planned? \\
Negative case & What named case would show that the upper-coherence reading is
unnecessary or false? \\
Demotion rule & What weaker sentence remains if the strong reading fails? \\
\bottomrule
\end{tabular}
\end{center}

This checklist is part of the public claim boundary. It protects the manuscript
from treating an attractive upper-level picture as completed science. It also
gives future work a precise route: every stronger \(K_{12}\) sentence must
bring the missing proof, finite witness, comparator, or negative case with it.

\subsubsection{Boundary and Falsifiability}

The candidate boundary \(\partial\Omega(K_{12})\) is crossed whenever the
declared upper-coherence reading loses one of its required witnesses. Examples
include recursive divergence, incompatible projections, operator assumptions
that cannot be stated, a missing demotion rule, or a lower-level witness that
changes the verdict after projection.

The candidate is falsified or demoted if:
\begin{itemize}
\item a proposed upper-level axis can be removed without changing any witness;
\item a claimed convergence criterion fails on a named counterexample;
\item an operator-symmetry statement is used without explicit operators;
\item a branch- or cycle-stability claim lacks a proof route or evidence row;
\item a lower-level domain example is used as universal evidence.
\end{itemize}

\subsubsection{Research Tests}

Future work may test explicitly modeled recursive meta-hierarchies, candidate
fixed-point behavior, operator assumptions, and branching constraints. Such
tests should be reported as research tests until they include source snapshots,
formal assumptions, negative controls, and reopening conditions. The present
release therefore treats \(K_{12}\) as a disciplined frontier object: useful
for organizing upper-bound questions, but not a current proof of global
closure, terminality, or universal invariance.

\subsubsection{Reader Guidance for the Upper Taxonomy}

A reader should use the \(K_{12}\) module as a map of questions rather than as
an answer sheet. The module asks what it would mean for a family of
meta-frameworks to become mutually inspectable, what kind of embedding context
would be required, which lower-level witnesses must remain visible, and which
operator assumptions would have to be stated before a stronger result could be
claimed. These questions are useful even while the answers remain provisional.

The safest reading is comparative. If an existing metatheory, formal ontology,
category-theoretic construction, type-theoretic universe, or scientific
unification program explains the same upper-level relation with a cleaner
support structure, the OC wording should become narrower. The K12 vocabulary
then remains a translation aid or a research prompt, not a priority claim and
not a proof of unique adequacy.

The module also protects the lower levels. Without an explicit upper-bound
frontier, broad synthesis language tends to leak backward into the formal core
and make \(K_0\)--\(K_{10}\) sound stronger than their evidence permits. By
keeping \(K_{12}\) provisional, the release can state a future direction
without allowing that direction to inflate current theorem, replay, or
comparator claims.

For practical readers, this means that \(K_{12}\) should not be used as an
enterprise, engineering, or scientific decision rule. It may help structure
questions about governance, evidence integration, and cross-domain
compatibility, but any operational recommendation must be grounded in a lower
level, a named source row, a comparator, a negative case, and a reopening
condition. The upper taxonomy is a research boundary, not an authorization
surface.

For formal readers, the frontier status is equally important. A future proof
would need to name the exact object, the universe or meta-space in which it
lives, the morphisms or operators allowed, the fixed or non-fixed points under
review, and the demotion condition for every retained witness. Until that
work is present, \(K_{12}\) is a disciplined conjectural scaffold. It is useful
because it says what would have to be proved, not because it has already been
proved.
